% !TeX program = xelatex
\documentclass[11pt,a4paper,UTF8,fontset=fandol]{ctexart}

% Keep this layout synchronized with the English course-report template.
\usepackage[margin=0.88in,headheight=14pt]{geometry}
\usepackage{graphicx}
\usepackage{booktabs}
\usepackage{amsmath,amssymb}
\usepackage{xcolor}
\usepackage{hyperref}
\usepackage{enumitem}
\usepackage{array}
\usepackage{needspace}
\usepackage{microtype}
\usepackage{titlesec}
\usepackage{fancyhdr}

% Shared assets live at the repository root.  The second path also supports
% compilers whose working directory is the repository root (for example Overleaf).
\graphicspath{{../assets/}{assets/}}

\hypersetup{
  colorlinks=true,
  linkcolor=blue!60!black,
  urlcolor=blue!60!black,
  citecolor=blue!60!black,
  pdfauthor={}
}

\setlength{\parskip}{0.45em}
\setlength{\parindent}{0pt}
\setlist[itemize]{topsep=0.25em,itemsep=0.15em,leftmargin=1.6em}
\setlist[enumerate]{topsep=0.25em,itemsep=0.15em,leftmargin=1.8em}

\titleformat{\section}{\Needspace{8\baselineskip}\normalfont\Large\bfseries}{\thesection}{0.8em}{}
\titleformat{\subsection}{\Needspace{6\baselineskip}\normalfont\large\bfseries}{\thesubsection}{0.8em}{}
\titlespacing*{\section}{0pt}{1.0em}{0.45em}
\titlespacing*{\subsection}{0pt}{0.75em}{0.35em}

% =========================
% 模板信息
% =========================
% 准备新报告时只需修改本区域。
\newcommand{\CourseName}{[课程名称]}
\newcommand{\ReportTitle}{[报告题目]}
\newcommand{\CoverReportTitle}{[报告题目第一行]\\[0.35cm]
[报告题目第二行]}
\newcommand{\StudentName}{[学生姓名]}
\newcommand{\StudentID}{[学生学号]}
\newcommand{\InstructorName}{[任课教师]}
\newcommand{\SchoolName}{[学院名称]}
\newcommand{\InstitutionName}{同济大学}
\newcommand{\ReportDate}{[年—月]}

% 首次只运行一次 XeLaTeX 时也显示一份简洁目录；生成 main.toc 后，
% 自动目录会接管并按正文更新页码。
\newcommand{\TemplateTocLine}[3][0pt]{%
  \par\noindent\hspace*{#1}#2\nobreak
  \leaders\hbox to 0.8em{\hss.\hss}\hfill\nobreak
  \makebox[2em][r]{#3}\par\vspace{0.35em}%
}
\newcommand{\TemplateTocPreview}{%
  \TemplateTocLine{1\quad 引言}{1}
  \TemplateTocLine{2\quad 问题陈述}{1}
  \TemplateTocLine{3\quad 方法}{1}
  \TemplateTocLine[2em]{3.1\quad 分析流程}{1}
  \TemplateTocLine[2em]{3.2\quad 实现说明}{1}
  \TemplateTocLine{4\quad 结果}{1}
  \TemplateTocLine{5\quad 讨论}{2}
  \TemplateTocLine{6\quad 结论}{2}
}
\makeatletter
\newcommand{\TemplateTableOfContents}{%
  \IfFileExists{\jobname.toc}{%
    \tableofcontents
  }{%
    \section*{目\quad 录}%
    \TemplateTocPreview
    \@starttoc{toc}%
  }%
}
\makeatother

% 页眉不显示课程、题目或个人信息。
\pagestyle{fancy}
\fancyhf{}
\fancyhead[L]{课程作业模板}
\fancyfoot[C]{\thepage}
\renewcommand{\headrulewidth}{0.3pt}

\begin{document}

% =========================
% 封面：结构与英文模板一致
% =========================
\begin{titlepage}
    \thispagestyle{empty}
    \newgeometry{left=2.0cm,right=2.0cm,top=1.8cm,bottom=2.0cm}

    \noindent
    \begin{minipage}[t]{0.36\textwidth}
        \vspace{0pt}
        \includegraphics[height=2.6cm]{top_left_logo.png}
    \end{minipage}
    \hfill
    \begin{minipage}[t]{0.58\textwidth}
        \vspace{0pt}
        \raggedleft
        \parbox[c][2.6cm][c]{\linewidth}{
            \raggedleft
            {\large\bfseries \InstitutionName 研究生课程\par}
            \vspace{0.35cm}
            {\large\bfseries \CourseName\par}
        }
    \end{minipage}

    \vspace{1.9cm}

    \begin{center}
        \includegraphics[width=0.58\textwidth]{tongji_logo.png}
    \end{center}

    \vspace{1.55cm}

    \begin{center}
        {\Huge\bfseries \CoverReportTitle\par}
        \vspace{1.05cm}
        {\LARGE\bfseries \CourseName\par}
    \end{center}

    \vspace{1.5cm}

    \begin{center}
        \renewcommand{\arraystretch}{2.0}
        \begin{tabular}{>{\Large\bfseries}r @{\hspace{0.8cm}} c}
            姓名： & \underline{\makebox[8.0cm][c]{\Large \StudentName}} \\
            学号： & \underline{\makebox[8.0cm][c]{\Large \StudentID}} \\
            任课教师： & \underline{\makebox[8.0cm][c]{\Large \InstructorName}} \\
        \end{tabular}
    \end{center}

    \vfill

    \begin{center}
        {\LARGE\bfseries \SchoolName \\ \InstitutionName\par}
        \vspace{0.8cm}
        {\Large \ReportDate\par}
    \end{center}

    \restoregeometry
\end{titlepage}

\pagenumbering{roman}
\TemplateTableOfContents
\newpage
\pagenumbering{arabic}

\section{引言}
本文档是一份中性的课程报告模板。请将本段替换为对选题、实践或理论动机以及报告所研究问题的简明说明。清晰的引言通常从较宽的背景逐步收束到边界明确的任务，并在结尾简要介绍全文结构。

模板不包含原作业的正文、数据、代码或插图，仅用于展示标题层级、公式、表格、列表、页码以及保留的封面版式。

\section{问题陈述}
在介绍方法之前，应先定义研究对象、基本假设、变量和评价标准。设 $x_t$ 表示系统状态，$u_t$ 表示决策变量，$c_t$ 表示观测条件，则一个通用目标函数可以写为
\begin{equation}
  J(\boldsymbol{u})=
  \sum_{t=1}^{T} L(x_t,u_t,c_t)+\lambda R(\boldsymbol{u}),
\end{equation}
其中 $L(\cdot)$ 衡量主要任务表现，$R(\cdot)$ 表示正则项或约束惩罚，$\lambda$ 用于控制两者之间的权衡。

\section{方法}
\subsection{分析流程}
一个清晰的方法部分可以按以下步骤组织：
\begin{enumerate}
  \item 说明输入、输出与基本假设；
  \item 描述模型或分析过程；
  \item 给出对比方法与评价指标；
  \item 解释如何处理不确定性与局限性。
\end{enumerate}

\subsection{实现说明}
本小节可用于记录理解报告所必需的实现细节。请避免写入私人文件路径、账户名称或设备专属信息。若未来提交需要代码或数据，可以说明其可用性，但不应嵌入个人标识信息。

\section{结果}
表~\ref{tab:example} 展示了一种紧凑的对比格式。正式使用时，请将所有标签和数值替换为课程任务中的真实结果。

\begin{table}[htbp]
  \centering
  \caption{示例对比结果表}
  \label{tab:example}
  \begin{tabular}{lccc}
    \toprule
    方法 & 指标 A & 指标 B & 综合评价 \\
    \midrule
    基准方法 & 0.72 & 0.68 & 参考 \\
    改进方法 & 0.81 & 0.76 & 提升 \\
    \bottomrule
  \end{tabular}
\end{table}

结果部分应先报告证据，再对其进行解释。应说明测量对象，并在适用时给出单位和不确定性，同时区分观察结果与原因分析。

\section{讨论}
讨论部分应解释观察到的现象为何合理、结果是否符合预期，以及分析可能在哪些条件下失效。重要局限可能包括假设简化、样本有限、测量不确定性或工况覆盖不足。尚未检验的判断应明确标为假设或未来工作。

\section{结论}
结论应回答引言中提出的问题。一个有力的结论会概括最可靠的发现、实际意义和分析边界，同时避免引入新的证据。

\end{document}
